\documentclass[11pt,letterpaper]{article}

\usepackage[top=1in, bottom=1in, left=1in, right=1in, headheight=14pt]{geometry}
\usepackage{fontspec}
\setmainfont{DejaVu Serif}
\setsansfont{DejaVu Sans}
\setmonofont{DejaVu Sans Mono}

\usepackage{xcolor}
\usepackage{titlesec}
\usepackage{fancyhdr}
\usepackage{enumitem}
\usepackage{hyperref}
\usepackage{booktabs}
\usepackage{tabularx}
\usepackage{longtable}
\usepackage{colortbl}
\usepackage{multirow}
\usepackage{array}
\usepackage{parskip}
\usepackage{tcolorbox}
\usepackage{graphicx}
\usepackage{lastpage}

%% ─── COLOR SCHEME: Philosophical/Mathematical ─────────────────────────────
\definecolor{primary}{HTML}{1A237E}       % Deep indigo — mathematical, cosmic
\definecolor{secondary}{HTML}{00695C}     % Teal — organic growth, nature
\definecolor{tertiary}{HTML}{E65100}      % Burnt orange — the fox, illumination
\definecolor{accent}{HTML}{283593}        % Rich indigo — highlights, arrows
\definecolor{lightprimary}{HTML}{E8EAF6}
\definecolor{lightsecondary}{HTML}{E0F2F1}
\definecolor{lighttertiary}{HTML}{FFF3E0}
\definecolor{rowgray}{HTML}{F5F5F5}
\definecolor{bordergray}{HTML}{BDBDBD}
\definecolor{subtlegray}{HTML}{757575}
\definecolor{darktext}{HTML}{212121}
\definecolor{goldaccent}{HTML}{F9A825}

%% ─── SECTION FORMATTING ───────────────────────────────────────────────────
\titleformat{\section}
  {\Large\bfseries\sffamily\color{primary}}
  {\thesection}{1em}{}
  [\vspace{-0.5em}\textcolor{bordergray}{\rule{\textwidth}{0.4pt}}]

\titleformat{\subsection}
  {\large\bfseries\sffamily\color{secondary}}
  {\thesubsection}{1em}{}

\titleformat{\subsubsection}
  {\normalsize\bfseries\sffamily\color{tertiary}}
  {\thesubsubsection}{1em}{}

\titlespacing*{\section}{0pt}{1.5em}{0.8em}
\titlespacing*{\subsection}{0pt}{1.2em}{0.5em}
\titlespacing*{\subsubsection}{0pt}{0.8em}{0.3em}

%% ─── CUSTOM ENVIRONMENTS ──────────────────────────────────────────────────
\newcommand{\stageheader}[2]{%
  \noindent\colorbox{#2}{\makebox[\textwidth][l]{%
    \textcolor{white}{\bfseries\sffamily\hspace{6pt}#1\hspace{6pt}}}}%
  \vspace{0.5em}\par%
}

\newtcolorbox{asciibox}{
  colback=rowgray, colframe=bordergray,
  arc=1pt, boxrule=0.5pt,
  left=6pt, right=6pt, top=4pt, bottom=4pt,
  fontupper=\small\ttfamily,
  breakable
}

\newtcolorbox{quotebox}{
  colback=lightprimary, colframe=primary,
  arc=2pt, boxrule=0.5pt,
  left=12pt, right=12pt, top=8pt, bottom=8pt,
  breakable
}

\newtcolorbox{tealbox}{
  colback=lightsecondary, colframe=secondary,
  arc=2pt, boxrule=0.5pt,
  left=12pt, right=12pt, top=8pt, bottom=8pt,
  breakable
}

\newtcolorbox{orangebox}{
  colback=lighttertiary, colframe=tertiary,
  arc=2pt, boxrule=0.5pt,
  left=12pt, right=12pt, top=8pt, bottom=8pt,
  breakable
}

\newcommand{\metafield}[2]{\textbf{\sffamily #1} #2\par}

%% ─── TABLE HELPERS ────────────────────────────────────────────────────────
\newcolumntype{L}[1]{>{\raggedright\arraybackslash}p{#1}}
\newcolumntype{C}[1]{>{\centering\arraybackslash}p{#1}}
\newcommand{\tableheader}[1]{\textbf{\sffamily\textcolor{white}{#1}}}

%% ─── HEADERS / FOOTERS ────────────────────────────────────────────────────
\pagestyle{fancy}
\fancyhf{}
\fancyhead[L]{\small\sffamily\textcolor{subtlegray}{College of Knowledge \textbar{} Complex Plane of Experience}}
\fancyhead[R]{\small\sffamily\textcolor{subtlegray}{GSD Mission Package}}
\fancyfoot[C]{\small\sffamily\textcolor{subtlegray}{Page \thepage\ of \pageref{LastPage}}}
\renewcommand{\headrulewidth}{0.4pt}
\renewcommand{\footrulewidth}{0pt}

%% ─── HYPERLINKS ───────────────────────────────────────────────────────────
\hypersetup{
  colorlinks=true,
  linkcolor=secondary,
  urlcolor=secondary,
  pdfauthor={GSD Ecosystem / Tibsfox},
  pdftitle={College of Knowledge and the Complex Plane of Experience -- Mission Package},
  pdfsubject={Vision to Mission Pipeline},
}

%% ══════════════════════════════════════════════════════════════════════════
\begin{document}

%% ─── TITLE PAGE ───────────────────────────────────────────────────────────
\thispagestyle{empty}
\begin{center}
\vspace*{2.5cm}

{\small\sffamily\textcolor{primary}{\textbf{GSD RESEARCH MISSION PACKAGE}}}

\vspace{0.3cm}
\textcolor{bordergray}{\rule{8cm}{0.4pt}}

\vspace{1.5cm}
{\fontsize{26}{32}\selectfont\bfseries\sffamily\textcolor{primary}{The College of Knowledge\\[0.3em]and the Complex Plane\\[0.3em]of Experience}}

\vspace{0.8cm}
{\Large\sffamily\textcolor{secondary}{Deep Linking, Seed Growth, and the Art of the Space Between}}

\vspace{0.5cm}
{\large\sffamily\itshape\textcolor{tertiary}{A Deep Retrospective and Lessons Learned}}

\vspace{2.5cm}
{\sffamily\textcolor{accent}{Vision $\rightarrow$ Research $\rightarrow$ Mission}}\\[0.3em]
{\small\sffamily\textcolor{accent}{Full Pipeline Execution --- Fleet Profile}}

\vspace{2.5cm}
{\small\sffamily\textcolor{subtlegray}{Date: March 26, 2026 \quad|\quad Status: Ready for GSD Orchestrator}}\\[0.3em]
{\small\sffamily\textcolor{subtlegray}{Activation Profile: Fleet (6 Modules) \quad|\quad Estimated: 18--22 context windows across 6--8 sessions}}\\[0.3em]
{\small\sffamily\textcolor{subtlegray}{Ecosystem: \texttt{github.com/Tibsfox/gsd-skill-creator} (dev branch)}}
\end{center}
\newpage

%% ─── TABLE OF CONTENTS ────────────────────────────────────────────────────
\tableofcontents
\newpage

%%═══════════════════════════════════════════════════════════════════════════
%% STAGE 1 — VISION DOCUMENT
%%═══════════════════════════════════════════════════════════════════════════
\stageheader{STAGE 1 --- VISION DOCUMENT}{primary}

\section{The College of Knowledge and the Complex Plane of Experience --- Vision Guide}

\metafield{Date:}{March 26, 2026}
\metafield{Status:}{Deep Retrospective --- Ready for Fleet Execution}
\metafield{Depends on:}{\texttt{gsd-skill-creator-analysis.md}, \texttt{cooking-with-claude-compiled-session.md},
\texttt{unit-circle-skill-creator-synthesis.md}, \texttt{gsd-mathematical-foundations-conversation.md},
\textit{The Space Between} (923 pages), \texttt{10-perl-cpan-addendum.md}}
\metafield{Context:}{This mission unifies three previously separate but deeply linked bodies of work ---
the College Structure architectural spec, the mathematical framework of the Complex Plane of Experience,
and the fractal growth pattern that connects them both. It is simultaneously a technical specification,
a pedagogical framework, and a philosophical retrospective.}

\subsection{Vision}

There is a moment in every great architectural discovery when you realize the structure you were
building already existed --- that you were not inventing, but uncovering. The College of Knowledge
arrived this way. It wasn't designed from first principles in a whiteboard session. It grew,
rhizomatically, from a series of conversations about cooking, mathematics, Perl, and the Amiga.
And in growing, it revealed something: the \textit{space between} knowledge domains is not empty.
It is the most structurally rich region of all.

The Complex Plane of Experience is the mathematical home of this discovery. In \textit{The Space
Between}, the plane is defined with two axes: Logic$\leftrightarrow$Creativity on the real axis,
Real$\leftrightarrow$Imaginary on the imaginary axis. Together, they produce a coordinate system
for all human understanding --- not as metaphor, but as working geometry. Every concept has a
position. Every skill has a trajectory. Every learning path traces an arc across the plane. The
unit circle is not an illustration of this model. It \textit{is} the model, made operational.

The College Structure is what happens when you build an institution on this plane rather than a
flat grid. Departments are not filing cabinets of categorized information. They are \textit{seed
organisms} --- minimal axioms that, faithfully iterated through Rosetta panels and calibration
cycles, grow into forests of interconnected understanding. A single concept --- ``exponential decay''
--- when planted in the Mathematics Department and expressed across seven language panels, generates
connections to cooking thermodynamics, signal processing, population biology, and the physics of
radioactive materials. The seed does not contain the forest. The grammar that grows the seed does.

This is the Amiga Principle applied to epistemology: remarkable outcomes through architectural
elegance and specialized execution paths, not brute-force content accumulation.

The deep linking problem is the hardest and most interesting part. A student who understands
Newton's cooling law in the Cooking Department should find that understanding \textit{already present}
when she arrives at the Calculus wing of the Mathematics Department. The knowledge shouldn't move ---
the student moves. The College's job is to ensure the landscape she walks through is coherent.
Cross-department chords (in the unit circle sense) make this possible: direct connections between
circle positions that would otherwise require traveling the full arc.

This mission retrospects what was built, surfaces the fractal patterns that emerged, formalizes
the deep linking architecture, and prepares a Fleet execution plan for building the complete
College seeding system.

\subsection{Problem Statement}

\begin{enumerate}
  \item \textbf{The flat-grid problem.} The original Knowledge Graph (16$\times$16 Cartesian) placed
  concepts at fixed positions with Euclidean distances. This was correct as a map but wrong as a
  learning model. Learning is not distance-minimization. It's angular traversal across the Complex
  Plane. Without the plane geometry, prerequisite chains are guesswork.

  \item \textbf{The orphaned-island problem.} Departments built in isolation become content silos.
  A Mathematics Department and a Cooking Department that share Newton's cooling law but don't know
  it means a student encounters the same concept twice, from different angles, with no bridge between.
  Deep links --- cross-department concept pointers --- don't yet exist as a first-class architectural
  feature.

  \item \textbf{The dormant-seed problem.} The College seeds (Mathematics, Cooking, etc.) exist in
  specification but have not been formally seeded into the skill-creator runtime. The fractal
  expansion machinery (Complex Plane positioning, calibration feedback, Rosetta panel routing)
  exists in architectural documents but not as executable code that produces living departments.

  \item \textbf{The retrospective debt.} Sixty-five milestones and 500k+ lines of code represent
  enormous architectural intelligence that has never been systematically documented as lessons
  learned. Patterns that were discovered through friction --- the DACP protocol, the staging layer,
  the Teacher/Student/Support model --- exist as implementations without formal retrospectives.
  Future architects will rediscover the same friction points.

  \item \textbf{The fractal-depth problem.} Fractal expansion promises that a simple seed self-replicates
  outward with increasing detail. But the depth limit is undefined. Without formal angular velocity
  constraints (from the unit circle model), a seed can over-expand into incoherence, or fail to
  expand past the first iteration because the calibration engine lacks context to guide growth.
\end{enumerate}

\subsection{Core Concept}

\begin{tealbox}
\textbf{The College of Knowledge is an L-system seeded on the Complex Plane of Experience.}
Each Department is an axiom. Each Rosetta Panel is a production rule. Each calibration cycle
is a generation. The output is not stored knowledge --- it is a living structure that grows
toward the specific learner standing at a specific angle on the unit circle.
\end{tealbox}

\subsubsection{Domain Map: The Complex Plane of Experience}

\begin{asciibox}
                    IMAGINARY AXIS
                   (Real Experience)
                          |
    CONCRETE              |              CONCRETE
    + CREATIVE           |             + LOGICAL
    (Art, cooking,       |           (Engineering,
     music, craft)      |            code, math)
                        |
CREATIVE  ──────────────+──────────────  LOGICAL
AXIS                    |              (Real Axis)
                        |
    ABSTRACT            |              ABSTRACT
    + CREATIVE          |             + LOGICAL
    (Philosophy,        |           (Theory, proofs,
     metaphor)         |            pure math)
                        |
                   (Imaginary Experience)
                   IMAGINARY AXIS

  Every skill has a (θ, r) position:
    θ = angle on the plane (abstract↔concrete, logic↔creativity)
    r = radius (maturity / depth of coverage)
    Tangent line at θ = activation interface for real-world application
\end{asciibox}

\subsubsection{Module Architecture}

\begin{asciibox}
  COLLEGE OF KNOWLEDGE — MODULE ARCHITECTURE
  ═══════════════════════════════════════════

  Module 1: COMPLEX PLANE FOUNDATIONS
  ├── Axes, coordinate system, angular positions
  ├── Unit circle as activation geometry
  └── Euler's formula as skill composition law

  Module 2: COLLEGE ARCHITECTURE
  ├── Department/Wing/Concept hierarchy
  ├── Rosetta panels (7-language expression)
  └── .college/ directory structure

  Module 3: DEEP LINKING ENGINE
  ├── Cross-department concept pointers
  ├── Chord detection (shortcut identification)
  └── Versine gap analysis (missing bridges)

  Module 4: SEED GROWTH & FRACTAL EXPANSION
  ├── L-system production rules for concept growth
  ├── Calibration engine as iteration driver
  └── Angular velocity constraints

  Module 5: RETROSPECTIVE & LESSONS LEARNED
  ├── 65-milestone audit (Amiga Principle instances)
  ├── Friction points → architectural patterns
  └── Teacher/Student/Support model origins

  Module 6: THE ART OF THE SPACE BETWEEN
  ├── Philosophical synthesis across authored works
  ├── Through-line: from quark to College
  └── What the gaps between departments contain
\end{asciibox}

\subsection{Research Modules}

\subsubsection{Module 1: Complex Plane Foundations}

The mathematical backbone of the entire system. This module formalizes the Complex Plane of
Experience as a working coordinate system: the two axes, the unit circle as activation geometry,
the tangent line as the skill-to-application interface, and Euler's formula ($e^{i\theta} = \cos\theta
+ i\sin\theta$) as the composition law for combining skills. This is the mathematical bridge
between \textit{The Space Between} (philosophical framework) and skill-creator (executable
system). The old 16$\times$16 Knowledge Graph is formally retired in favor of $(θ, r)$ polar
coordinates on the plane.

\subsubsection{Module 2: College Architecture}

The `.college/` directory structure, its Department/Wing/Concept hierarchy, the Rosetta panels
(7-language expression: Python, C++, Java, Perl, Lisp, Pascal, Fortran), and the three-tier
knowledge organization (pedagogical annotation + panel expression + calibration profile). This
module maps the CPAN architectural lineage (PAUSE $\to$ contribution pipeline, CPAN Index $\to$
concept registry, BackPAN $\to$ delta store) directly onto the College Structure's design
constraints. It also specifies the Mathematics Department seed in full detail as the reference
implementation.

\subsubsection{Module 3: Deep Linking Engine}

The hardest unsolved problem: how does a concept in one Department ``know about'' its equivalent
in another? This module formalizes cross-department concept pointers as \textit{chords} on the
unit circle --- direct connections between angular positions that bypass the arc path. Chord
detection is an algorithmic problem (given two concept $(θ, r)$ positions, compute the straight-line
distance and tag as a cross-department bridge if shorter than arc travel). Versine gap analysis
identifies \textit{missing} links: where a bridge should exist (two concepts near each other
on the plane) but doesn't yet.

\subsubsection{Module 4: Seed Growth and Fractal Expansion}

The L-system formalization of how a minimal Department seed expands into a full knowledge
ecosystem. Production rules govern how a parent concept generates child concepts; angular velocity
constraints prevent over-expansion; calibration feedback (Observe $\to$ Compare $\to$ Adjust $\to$
Record) drives each generation. This module defines the mathematical limits of fractal depth ---
the conditions under which expansion produces coherence vs.\ noise --- and specifies the
Mandelbrot-equivalent boundary condition for College growth.

\subsubsection{Module 5: Retrospective and Lessons Learned}

A systematic audit of all 65+ milestones through the lens of the Amiga Principle. This module
extracts the \textit{friction-to-pattern} transformations: every place where the build hit a
wall, and what architectural insight the wall revealed. Key discoveries to document include the
DACP protocol origin (markdown ambiguity $\to$ three-part bundle), the staging layer emergence,
the Teacher/Student/Support model, the planning-bridge MCP server, and the virtual Black Rock
City experiment. This is the ecosystem's institutional memory.

\subsubsection{Module 6: The Art of the Space Between}

The philosophical synthesis. This module weaves the through-line across all five authored works
(\textit{The Space Between}, \textit{The Hundred Voices}, \textit{The Quark and the Compiler},
\textit{The Longer Road}, and the GSD ecosystem itself as a sixth ``living work''). It asks: what
\textit{is} the space between departments in the College? What lives in the chords, the versine
gaps, the tangent extensions? The answer, consistent across every scale of the ecosystem, is that
the space between is not absence. It is the highest-density region of meaning --- the morphism in
category theory, the functor that preserves structure across maps, the Coast Salish concept of
\textit{everything in relation}.

\subsection{Chipset Configuration}

\begin{asciibox}
# chipset.yaml — College of Knowledge Mission

chipset:
  name: "college-knowledge-complex-plane"
  version: "1.0.0"
  profile: fleet

  agnus:           # DMA / orchestration
    role: FLIGHT
    model: claude-opus-4-20250514
    context_windows: 4
    responsibilities:
      - Module synthesis and integration
      - Cross-department chord validation
      - Philosophical through-line coherence

  denise:          # Display / pattern generation
    role: CRAFT-FRACTAL
    model: claude-opus-4-20250514
    context_windows: 3
    responsibilities:
      - L-system production rule formalization
      - Fractal depth boundary conditions
      - Angular velocity constraint derivation

  paula:           # Audio / calibration (feedback loops)
    role: CRAFT-CALIBRATION
    model: claude-sonnet-4-20250514
    context_windows: 3
    responsibilities:
      - Calibration engine specification
      - Observe-Compare-Adjust-Record formalization
      - Feedback loop integration testing

  68000:           # CPU / execution
    role: EXEC (×2)
    model: claude-sonnet-4-20250514
    context_windows: 6
    responsibilities:
      - Document production (parallel tracks)
      - Rosetta panel expressions
      - College seed implementations

  support:
    role: HAIKU scaffolding
    model: claude-haiku-4-5-20251001
    context_windows: 2
    responsibilities:
      - Schema generation
      - Template scaffolding
      - Wave 0 foundation
\end{asciibox}

\subsection{Scope Boundaries}

\subsubsection{In Scope (v1.0)}
\begin{itemize}[noitemsep]
  \item Full formalization of the Complex Plane of Experience as executable coordinate system
  \item `.college/` directory structure specification and reference seed implementation (Mathematics)
  \item Deep linking engine: chord detection algorithm, versine gap analysis, cross-department pointer format
  \item L-system production rules for Department seed growth, angular velocity constraints
  \item 65-milestone retrospective: friction points, architectural patterns, institutional memory document
  \item Philosophical synthesis: through-line across all authored works
  \item Test suite for plane positioning, chord computation, and seed expansion
\end{itemize}

\subsubsection{Out of Scope}
\begin{itemize}[noitemsep]
  \item Full implementation of all College Departments (Mathematics is the reference; others are seeded v2.0+)
  \item Perl panel implementation (addressed in \texttt{10-perl-cpan-addendum.md})
  \item Wasteland/DoltHub federation integration (separate mission)
  \item 3D extension of the Complex Plane (v3.0 consideration)
  \item Public documentation site update (frozen pending architecture finalization)
\end{itemize}

\subsection{Success Criteria}

\begin{enumerate}
  \item The Complex Plane coordinate system is implemented and returns a $(θ, r)$ position for any
  given concept with deterministic consistency ($\pm 0.01$ radians tolerance).

  \item The `.college/` directory structure exists with complete schema documentation; the Mathematics
  Department seed compiles and loads without errors.

  \item The Rosetta Core routes any concept query to the correct panel based on its plane position
  (systems/heritage/bridge routing verified against 20 test concepts).

  \item The chord detection algorithm identifies cross-department bridges with $\geq 90\%$ precision
  against a hand-labeled test set of 30 known concept pairs.

  \item The versine gap analysis produces a ranked list of missing cross-department bridges, with
  the top 10 verified as legitimate gaps by human review.

  \item L-system production rules expand the Mathematics seed to at least 3 generations without
  coherence degradation (calibration engine confirms).

  \item The angular velocity constraint system prevents over-expansion: a bounded seed stays within
  defined plane region with $> 95\%$ of generated concepts remaining coherent.

  \item The 65-milestone retrospective document is complete: every milestone has at least one
  friction point and one architectural pattern derived from it.

  \item The friction-to-pattern registry contains at least 15 distinct patterns with formal names,
  descriptions, and instances.

  \item All authored works (\textit{The Space Between}, \textit{The Hundred Voices}, \textit{The
  Quark and the Compiler}, \textit{The Longer Road}) have their through-line mapped to specific
  College Architecture decisions.

  \item The deep linking test suite passes: 95\%+ of cross-department navigation paths resolve
  correctly from any starting Department to any target concept.

  \item The complete package is self-contained: a new engineer can build and run the College seed
  system from this document alone, with no additional context.
\end{enumerate}

\subsection{Relationship to Other Documents}

\begin{center}
\rowcolors{2}{rowgray}{white}
\begin{tabularx}{\textwidth}{L{4.5cm}L{2.2cm}X}
\toprule
\rowcolor{primary}
\tableheader{Document} & \tableheader{Relationship} & \tableheader{What It Contributes} \\
\midrule
\textit{The Space Between} (923pp) & Foundation & Complex Plane of Experience framework, set-theoretic identity model, fractal geometry \\
\texttt{unit-circle-skill-creator-synthesis.md} & Prerequisite & Unit circle as activation geometry, tangent/chord/versine mechanics \\
\texttt{cooking-with-claude-compiled-session.md} & Parallel & College Structure spec, Rosetta Core, calibration engine, Mathematics seed \\
\texttt{gsd-mathematical-foundations-conversation.md} & Prerequisite & L-system formalization, category theory as morphism model \\
\texttt{10-perl-cpan-addendum.md} & Extension & Perl panel, CPAN $\to$ College lineage, 7th panel routing \\
\texttt{gsd-skill-creator-analysis.md} & Context & Current codebase state (500k+ LOC, 65 milestones) \\
\texttt{gsd-bbs-educational-pack-vision.md} & Parallel & BBS as historical College (FidoNet $\to$ Kafka functor) \\
\texttt{gsd-chipset-architecture-vision.md} & Context & Amiga chip naming convention, specialized execution paths \\
\bottomrule
\end{tabularx}
\end{center}

\subsection{The Through-Line}

The Amiga was not the most powerful computer of its era. It was the most architecturally leveraged.
Agnus handled DMA so the CPU could think. Denise handled display so the CPU could compute. Paula
handled audio so the CPU could reason about structure rather than samples. The chips did not share
work randomly --- each had a domain, a specialty, a piece of the whole that only it could hold.

The College of Knowledge is an Amiga.

Each Department is a chip. Mathematics handles the DMA of abstraction --- the background movement
of formal structure through the system. Cooking handles the display layer --- the sensory, concrete
interface where understanding becomes visible and delicious. Music handles the audio layer --- the
temporal, rhythmic, emotional channel. Each Department expresses the same underlying concepts
(energy, transformation, pattern, composition) through its specialized substrate, and in doing so
frees the learner's CPU to do what only humans can: \textit{stand in the space between and see
what the chips cannot.}

The space between departments is not a gap to be filled. It is the highest-value real estate in
the entire system. It is where Euler lives --- that astonishing moment when $e^{i\pi} + 1 = 0$,
and you realize that five apparently unrelated mathematical constants are secretly one conversation.
It is where the cook who understands heat transfer meets the engineer who understands differential
equations, and they discover they've been solving the same problem in different vocabularies.

The art of the space between is not decoration. It is the curriculum.

\newpage

%%═══════════════════════════════════════════════════════════════════════════
%% STAGE 2 — RESEARCH REFERENCE
%%═══════════════════════════════════════════════════════════════════════════
\stageheader{STAGE 2 --- RESEARCH REFERENCE}{secondary}

\section{College of Knowledge and Complex Plane --- Research Reference}

\subsection{How to Use This Document}

This research reference synthesizes findings across the GSD ecosystem's primary sources. Each
module's findings are organized as self-contained reference material for the corresponding
execution wave. Sources are from authored works, architectural conversation transcripts, and
peer-validated mathematical frameworks.

\textbf{Key source authorities:}
\begin{itemize}[noitemsep]
  \item \textit{The Space Between} (Tibsfox / Miles Tiberius Foxglove) --- primary philosophical and mathematical source
  \item \texttt{unit-circle-skill-creator-synthesis.md} --- unit circle geometry formalization
  \item \texttt{cooking-with-claude-compiled-session.md} --- College Structure specification
  \item \texttt{gsd-mathematical-foundations-conversation.md} --- L-system and category theory framework
  \item \texttt{10-perl-cpan-addendum.md} --- CPAN architectural lineage
  \item \textit{The Hundred Voices} (Tibsfox) --- fractal narrative structure across 100 authorial voices
\end{itemize}

\subsection{Module 1 Reference: Complex Plane Foundations}

\subsubsection{The Two Axes}

The Complex Plane of Experience as defined in \textit{The Space Between} establishes a coordinate
system for all human experience and understanding:

\begin{center}
\rowcolors{2}{rowgray}{white}
\begin{tabularx}{\textwidth}{L{3cm}L{4cm}L{4cm}X}
\toprule
\rowcolor{primary}
\tableheader{Axis} & \tableheader{Negative Pole} & \tableheader{Positive Pole} & \tableheader{What It Measures} \\
\midrule
Real (horizontal) & Pure Logic / Analysis & Pure Creativity / Synthesis & Balance between analytical and generative thinking \\
Imaginary (vertical) & Purely Abstract / Theoretical & Purely Concrete / Experiential & Distance from direct sensory/practical engagement \\
\bottomrule
\end{tabularx}
\end{center}

Every concept, skill, and learning path has a position on this plane. A pure mathematical proof
sits deep in the Logic + Abstract quadrant. A cooking technique sits in the Creativity + Concrete
quadrant. A recipe that teaches thermodynamics crosses quadrants --- and that crossing is precisely
why it works as pedagogy.

\subsubsection{Unit Circle as Activation Geometry}

The unit circle synthesis established that every classical trigonometric function maps to a
concrete skill-creator function:

\begin{center}
\rowcolors{2}{rowgray}{white}
\begin{tabularx}{\textwidth}{L{2cm}L{4cm}X}
\toprule
\rowcolor{secondary}
\tableheader{Segment} & \tableheader{Geometric Identity} & \tableheader{Skill-Creator Function} \\
\midrule
sin($\theta$) & Height of point on circle & Abstract reach of a skill --- how far into theory it extends \\
cos($\theta$) & Horizontal grounding & Concrete applicability --- how much is directly usable \\
tan($\theta$) & Reach beyond the circle & Activation reach --- how far the skill extends into real-world application \\
arc & Curved path along circumference & The learning path --- distance traveled around the plane \\
chord & Straight line between two positions & Cross-department shortcut --- the deep link \\
versine & $1 - \cos\theta$ (gap from full grounding) & Gap analysis --- what's missing from complete understanding \\
exsecant & sec$\theta - 1$ (reach beyond circle) & External resource reach --- MCP integrations, community skills \\
\bottomrule
\end{tabularx}
\end{center}

\textbf{Critical insight:} Point A (where the circle meets the tangent line) is the interface
between bounded abstract exploration and unbounded real-world application. The direction of
application is entirely determined by the angle of abstract exploration. Change $\theta$ and the
same body of abstract knowledge connects to completely different concrete applications.

\subsubsection{Euler's Formula as Composition Law}

$e^{i\theta} = \cos\theta + i\sin\theta$ is not merely a mathematical curiosity --- it is the
composition law for skill-creator. When two skills at angles $\theta_1$ and $\theta_2$ combine,
the result is their product as complex numbers: $e^{i(\theta_1 + \theta_2)}$. This means:

\begin{itemize}[noitemsep]
  \item Skills compose by \textit{adding angles} --- the combined skill's plane position is predictable
  \item The radius (depth/maturity) multiplies
  \item Skills at opposite angles ($\theta$ and $\theta + \pi$) cancel --- they are true complements
  \item The identity skill (no transformation) sits at angle $0$ --- it does not change the learner's position
\end{itemize}

\subsection{Module 2 Reference: College Architecture}

\subsubsection{The CPAN Lineage}

The College Structure did not emerge from a vacuum. Its direct intellectual ancestor is CPAN (1995),
which established the architectural patterns that have governed every major knowledge ecosystem since:

\begin{center}
\rowcolors{2}{rowgray}{white}
\begin{tabularx}{\textwidth}{L{4cm}L{4cm}X}
\toprule
\rowcolor{primary}
\tableheader{CPAN Component} & \tableheader{College Equivalent} & \tableheader{Principle} \\
\midrule
PAUSE upload server & Contribution Pipeline & Human-reviewed ingestion with namespace protection \\
CPAN Index & Concept Registry & Canonical mapping of names to locations \\
Module namespaces & Department/Wing/Concept hierarchy & Hierarchical knowledge organization \\
META.yml & RosettaConcept metadata & Structured metadata: dependencies, version, description \\
CPAN Testers & Calibration Engine gates & Automated cross-environment quality assurance \\
Mirror network & Progressive disclosure tiers & Distributed, cached copies at multiple fidelity levels \\
BackPAN (archive) & Delta Store & Complete history retained, nothing truly deleted \\
Developer releases & Observation pipeline (draft) & Tentative knowledge that doesn't pollute the index \\
\texttt{perldoc} viewer & College Loader's explore() & Browse knowledge without modifying it \\
POD embedded docs & Code-as-curriculum & Documentation lives with the code it describes \\
Dependency resolution & Concept prerequisite loading & Load what you need, resolve chains automatically \\
\bottomrule
\end{tabularx}
\end{center}

\subsubsection{Three-Tier Knowledge Organization}

The College Structure inherited the ``cookbook'' three-tier pattern from Numerical Recipes (Press et al.)
and refined it for executable knowledge:

\begin{center}
\rowcolors{2}{rowgray}{white}
\begin{tabularx}{\textwidth}{L{3.5cm}L{3.5cm}X}
\toprule
\rowcolor{secondary}
\tableheader{Tier} & \tableheader{Name} & \tableheader{Contents} \\
\midrule
Tier 1 & Pedagogical Annotation & Why this concept matters; connections to other concepts; historical context \\
Tier 2 & Panel Expressions & 7 Rosetta panel implementations (Python, C++, Java, Perl, Lisp, Pascal, Fortran) \\
Tier 3 & Calibration Profiles & Known good parameters; feedback-refined defaults; edge case handling \\
\bottomrule
\end{tabularx}
\end{center}

\subsubsection{The Seven Rosetta Panels}

\begin{center}
\rowcolors{2}{rowgray}{white}
\begin{tabularx}{\textwidth}{L{2.5cm}L{2.5cm}X}
\toprule
\rowcolor{primary}
\tableheader{Panel} & \tableheader{Family} & \tableheader{What It Teaches} \\
\midrule
Python & Systems & Readability, scientific computing, rapid prototyping \\
C++ & Systems & Memory management, performance, templates \\
Java & Systems & Type safety, OOP patterns, enterprise architecture \\
Perl & Bridge & Linguistics, regex as syntax, CPAN ecosystem architecture, TMTOWTDI \\
Lisp & Heritage & Functional programming, homoiconicity, code-as-data \\
Pascal & Heritage & Structured programming, educational clarity, Wirth's principles \\
Fortran & Heritage & Numerical computing, array operations, scientific heritage \\
\bottomrule
\end{tabularx}
\end{center}

The Perl ``bridge'' panel is uniquely positioned on the Complex Plane: it spans systems and heritage
families. Designed by a trained linguist (Larry Wall), Perl's syntax encodes Huffman coding
(common operations are concise), context sensitivity, and TMTOWTDI (``There's More Than One Way
To Do It'') --- the Rosetta principle expressed as language philosophy.

\subsection{Module 3 Reference: Deep Linking Engine}

\subsubsection{Chord Geometry}

A chord on the unit circle connects two points $(θ_1, r_1)$ and $(θ_2, r_2)$. Its length is:

\[
  \text{chord} = \sqrt{r_1^2 + r_2^2 - 2r_1 r_2 \cos(\theta_2 - \theta_1)}
\]

The arc distance between the same two points (staying on the circle at radius $r$) is:

\[
  \text{arc} = r \cdot |\theta_2 - \theta_1|
\]

A deep link is warranted when chord $\ll$ arc --- when the straight-line connection is significantly
shorter than the arc path. The threshold is empirically calibrated: a chord/arc ratio below $0.6$
is a strong deep link candidate.

\subsubsection{Known Cross-Department Bridges}

\begin{center}
\rowcolors{2}{rowgray}{white}
\begin{tabularx}{\textwidth}{L{3.5cm}L{3.5cm}L{2.5cm}X}
\toprule
\rowcolor{secondary}
\tableheader{Concept A (Dept)} & \tableheader{Concept B (Dept)} & \tableheader{Bridge Type} & \tableheader{Shared Structure} \\
\midrule
Newton's cooling (Math) & Maillard reaction temp (Cooking) & Exponential decay & $dT/dt = k(T - T_\text{env})$ \\
Fourier series (Math) & Timbre/overtones (Music) & Frequency decomposition & Harmonic analysis \\
Fibonacci sequence (Math) & Phyllotaxis (Biology) & Irrational growth ratio & Golden ratio $\phi$ \\
Gradient descent (Math) & Instrument tuning (Music) & Error minimization & Feedback correction \\
Entropy (Physics) & Sourdough fermentation (Cooking) & Thermodynamic disorder & Spontaneous state change \\
Set theory (Math) & Recipe ingredient sets (Cooking) & Membership functions & Inclusion/exclusion \\
L-systems (Math) & Branching recipes (Cooking) & Production grammar & Recursive substitution \\
\bottomrule
\end{tabularx}
\end{center}

\subsubsection{Versine Gap Analysis}

The versine ($1 - \cos\theta$) measures the gap between a concept's current grounding and complete
concretization. Applied to the College, a high versine score means a concept is theoretically
present but not yet concretely expressed. The gap analysis produces a ranked list of ``missing
bridges'' --- concept pairs whose plane positions are close (high cosine similarity) but whose
chord link does not yet exist in the College Structure.

\subsection{Module 4 Reference: Seed Growth and Fractal Expansion}

\subsubsection{The L-System Model}

From the mathematical foundations work, the GSD ecosystem is formally an L-system (Lindenmayer
system), a string-rewriting grammar originally developed to model plant growth:

\begin{center}
\rowcolors{2}{rowgray}{white}
\begin{tabularx}{\textwidth}{L{3.5cm}X}
\toprule
\rowcolor{primary}
\tableheader{L-System Component} & \tableheader{College Equivalent} \\
\midrule
Axiom & Minimal Department seed (e.g., ``complex numbers'') \\
Production rules & Rosetta panel expressions that generate child concepts \\
Iteration (generation $n$) & Each calibration cycle \\
Output string & The full concept graph of a Department at generation $n$ \\
Bracket notation $[\ ]$ & Branch protection --- concepts that grow but don't alter the parent \\
Stochastic variation & The human learner's unique derivation path \\
\bottomrule
\end{tabularx}
\end{center}

\textbf{Key property:} Some children will sprint across that bridge. Some will stop halfway and
build something beautiful where they are. Some will find a side path that nobody knew existed.
That's an L-system with stochastic variation. The grammar doesn't dictate the tree. It enables
the tree.

\subsubsection{Fractal Depth Constraints}

Without boundary conditions, L-system expansion produces infinite regress. The angular velocity
constraint limits how fast a seed's $\theta$ position can rotate per generation:

\[
  \left|\frac{d\theta}{dn}\right| \leq \omega_\text{max}
\]

where $\omega_\text{max}$ is calibrated per learner context. A highly concrete learner (low $r$,
concrete $\theta$) has a lower $\omega_\text{max}$ --- concepts should stay near the concrete axis.
An advanced learner exploring abstract connections has a higher $\omega_\text{max}$.

The Mandelbrot boundary condition for College growth: a concept remains coherent if, iterated
through the production rules, its orbit on the plane stays bounded. If it escapes to infinity
(incoherence), the calibration engine terminates that branch.

\subsubsection{The Calibration Loop as Generation Driver}

\begin{center}
\rowcolors{2}{rowgray}{white}
\begin{tabularx}{\textwidth}{L{2.5cm}L{2cm}X}
\toprule
\rowcolor{secondary}
\tableheader{Step} & \tableheader{L-System Role} & \tableheader{What Happens} \\
\midrule
Observe & Read string & Current concept graph state is captured \\
Compare & Apply production rules & New concepts are generated from current state \\
Adjust & Prune/promote & Coherent branches promoted; incoherent branches pruned \\
Record & Write string (generation $n+1$) & New concept graph state is stored \\
\bottomrule
\end{tabularx}
\end{center}

\subsection{Module 5 Reference: Retrospective and Lessons Learned}

\subsubsection{The Friction-to-Pattern Registry (Selected Entries)}

\begin{center}
\rowcolors{2}{rowgray}{white}
\begin{tabularx}{\textwidth}{L{3.5cm}L{4cm}X}
\toprule
\rowcolor{primary}
\tableheader{Friction Point} & \tableheader{Pattern Derived} & \tableheader{Where It Lives Now} \\
\midrule
Markdown ambiguity in agent handoffs & DACP: three-part bundle (intent + data + code) & v1.49, agent communication protocol \\
Agent overwhelm with full codebase context & Staging layer / inbox pattern & Planning bridge MCP server \\
Teacher explaining to student who knows more & Teacher/Student/Support model (Opus/Sonnet/Haiku) & v1.50 pedagogy design \\
Flat grid failing to model prerequisite chains & Complex Plane $(θ, r)$ positioning & College Architecture spec \\
Skills isolated by language & Rosetta Core multi-panel expression & Cooking with Claude mission \\
Single-purpose skills over-specialized & Calibration engine generalizes from feedback & Fractal expansion machinery \\
Documentation 25 milestones behind code & Institutional retrospective mission (this doc) & v1.50 Half A audit \\
Over-eager integration proposals & Comprehension-before-integration principle & Ecosystem working patterns \\
\bottomrule
\end{tabularx}
\end{center}

\subsubsection{The Amiga Principle in Practice: Milestone Instances}

Every time the ecosystem encountered a problem that brute-force (more tokens, larger models,
longer prompts) would solve expensively, the Amiga Principle found an architectural solution:

\begin{itemize}[noitemsep]
  \item \textbf{DACP (v1.49):} Instead of longer markdown prompts between agents, structured
  three-part bundles. Smaller. Faster. Deterministic. The Agnus chip solution.

  \item \textbf{Staging layer:} Instead of loading full codebase context every session, a
  planning-bridge MCP server stages work in \texttt{.planning/staging/inbox/}. The Denise chip
  solution --- offload to the display chip, free the CPU.

  \item \textbf{Teacher/Student/Support:} Instead of one Opus model doing everything, specialized
  roles with appropriate model sizing. 30\% Opus judgment, 60\% Sonnet structure, 10\% Haiku
  scaffolding. The Paula chip solution --- audio doesn't need a full CPU.

  \item \textbf{Rosetta panels:} Instead of one monolithic language implementation, seven
  specialized panel expressions that together cover every learner angle. Amiga custom chips,
  not generic processing.
\end{itemize}

\subsection{Module 6 Reference: The Art of the Space Between}

\subsubsection{Authored Works Through-Line}

\begin{center}
\rowcolors{2}{rowgray}{white}
\begin{tabularx}{\textwidth}{L{4.5cm}L{3cm}X}
\toprule
\rowcolor{tertiary}
\tableheader{Work} & \tableheader{Core Claim} & \tableheader{Connection to College Architecture} \\
\midrule
\textit{The Space Between} (923pp) & The maps between things matter more than the things & The morphism \textit{is} the curriculum; cross-department chords are where learning lives \\
\textit{The Hundred Voices} (733pp) & 100 authorial voices spanning 1926--2025; the forest, not the trees & Fractal expansion: same axiom, stochastic variation, every voice a unique derivation \\
\textit{The Quark and the Compiler} & Tracing one particle through 13.8B years of grammar & L-system trace: the quark is the seed, the production rules are physics, the output is consciousness \\
\textit{The Longer Road} & Incremental, grounded progress over rushing to completion & Angular velocity constraint: don't over-rotate; the path matters \\
GSD ecosystem (living work) & Giving people their lives back & The endgame: when the College handles friction, humans focus on family, art, music, connection \\
\bottomrule
\end{tabularx}
\end{center}

\subsubsection{Category Theory as the Deep Link Formalism}

Category theory provides the formal language for the space between. A functor is a structure-preserving
map between categories. The deep links in the College are functors: they don't just connect two
concepts, they preserve the \textit{structure of relationships} between them. When Newton's
cooling law in the Math Department connects to the Maillard reaction in the Cooking Department,
the functor doesn't just say ``these are related.'' It says: every derivative in the Math structure
maps to a rate of change in the Cooking structure. Every boundary condition maps to a culinary
threshold. The entire abstract structure is preserved across the map.

This is why \textit{The Space Between} is a category-theoretic claim as a title: what matters is
not what humans are or what machines are, but the morphism between them.

\subsubsection{Coast Salish Cultural Context}

The concept of the space between has deep roots in Coast Salish worldview: everything in relation,
the cedar and the salmon and the raven and the human as nodes in a network where the \textit{connections}
carry the meaning. The College Architecture embeds this principle structurally: no Department exists
in isolation; every concept carries its relationship metadata; the between-space is first-class.
Nation-specific attribution: this is specifically Coast Salish epistemology, not a generic
``Indigenous perspective.''

\subsection{Safety and Sensitivity Considerations}

\begin{center}
\rowcolors{2}{rowgray}{white}
\begin{tabularx}{\textwidth}{L{3.5cm}L{2cm}X}
\toprule
\rowcolor{primary}
\tableheader{Consideration} & \tableheader{Boundary} & \tableheader{Approach} \\
\midrule
Coast Salish cultural reference & GATE & Use specific nation attribution; never generic ``Indigenous''; validate with OCAP\textregistered{} principles \\
Learner data from calibration & GATE & Calibration profiles are private by default; shared data requires explicit opt-in \\
Mathematical claims & ANNOTATE & All quantitative claims (angular velocity, chord thresholds) are empirically calibrated, not derived from first principles \\
Authored work citations & ABSOLUTE & All citations to authored works use exact titles and pen name; never paraphrase core philosophical claims \\
\bottomrule
\end{tabularx}
\end{center}

\subsection{Source Bibliography}

\textbf{Primary Sources (Authored):}
\begin{itemize}[noitemsep]
  \item Foxglove, M.T. (Tibsfox). \textit{The Space Between: The Autodidact's Guide to the Galaxy.} 923 pp. Available: \url{https://tibsfox.com/media/The_Space_Between.pdf}
  \item Tibsfox. \textit{The Hundred Voices.} 733 pp., 100 authorial voices, 1926--2025.
  \item Tibsfox. \textit{The Quark and the Compiler.} Philosophical essay.
  \item Tibsfox. \textit{The Longer Road.} Philosophical essay.
\end{itemize}

\textbf{GSD Ecosystem Documents:}
\begin{itemize}[noitemsep]
  \item \texttt{unit-circle-skill-creator-synthesis.md} — Unit circle geometry formalization
  \item \texttt{cooking-with-claude-compiled-session.md} — College Structure specification and Rosetta Core
  \item \texttt{gsd-mathematical-foundations-conversation.md} — L-system, category theory, set theory frameworks
  \item \texttt{10-perl-cpan-addendum.md} — Perl panel, CPAN architectural lineage
  \item \texttt{gsd-skill-creator-analysis.md} — Current codebase architecture (500k+ LOC, 65 milestones)
\end{itemize}

\textbf{Mathematical Foundations:}
\begin{itemize}[noitemsep]
  \item Lindenmayer, A. (1968). Mathematical models for cellular interactions in development. \textit{Journal of Theoretical Biology}, 18(3), 280--315. [L-system original paper]
  \item Mandelbrot, B.B. (1982). \textit{The Fractal Geometry of Nature.} W.H. Freeman.
  \item Mac Lane, S. (1971). \textit{Categories for the Working Mathematician.} Springer. [Category theory / functors]
  \item Wall, L., Christiansen, T., \& Orwant, J. (2000). \textit{Programming Perl} (3rd ed.). O'Reilly.
  \item Press, W.H. et al. (1992). \textit{Numerical Recipes in C: The Art of Scientific Computing.} Cambridge University Press.
\end{itemize}

\newpage

%%═══════════════════════════════════════════════════════════════════════════
%% STAGE 3 — MISSION PACKAGE
%%═══════════════════════════════════════════════════════════════════════════
\stageheader{STAGE 3 --- MISSION PACKAGE}{tertiary}

\section{Milestone Specification}

\subsection{Mission Objective}

Formalize the College of Knowledge as an executable L-system seeded on the Complex Plane of
Experience, implement the deep linking engine for cross-department chord navigation, produce a
systematic retrospective of the 65+ milestone journey, and synthesize the philosophical
through-line across all authored works into a living architectural document. The mission operates
as a Fleet-profile execution across six parallel-capable modules, culminating in a complete
College seeding system ready for handoff to the GSD orchestrator.

\subsection{Architecture Overview}

\begin{asciibox}
  FLEET EXECUTION ARCHITECTURE
  ══════════════════════════════════════════════════════

  Wave 0 (Sequential, Haiku):
  └── Foundation schemas, concept metadata format, test fixtures

  Wave 1A (Parallel Track A, Sonnet):         Wave 1B (Parallel Track B, Opus):
  ├── Module 1: Complex Plane math impl.       ├── Module 5: Retrospective audit
  └── Module 2: College Architecture spec      └── Module 6: Philosophical synthesis

  Wave 1C (Parallel Track C, Sonnet):
  ├── Module 3: Deep Linking engine
  └── Module 4: Seed Growth / L-system

  Wave 2 (Sequential, Opus):
  └── Cross-module integration: plane ↔ links ↔ seeds ↔ retrospective

  Wave 3 (Sequential, Sonnet):
  └── Publication: compiled College doc + test suite + index.html

  ══════════════════════════════════════════════════════
  Dependencies: 1A,1B,1C must complete before Wave 2
                Wave 2 must complete before Wave 3
\end{asciibox}

\subsection{System Layers}

\begin{enumerate}[noitemsep]
  \item \textbf{Foundation Layer} (Wave 0) — Shared schemas, test fixtures, concept metadata format
  \item \textbf{Mathematical Layer} (Wave 1A) — Complex Plane coordinate system, unit circle geometry, Euler composition
  \item \textbf{Architectural Layer} (Wave 1A) — College directory structure, Department seeds, Rosetta panel routing
  \item \textbf{Linking Layer} (Wave 1C) — Chord detection, versine gap analysis, cross-department pointer format
  \item \textbf{Growth Layer} (Wave 1C) — L-system rules, angular velocity constraints, calibration loop
  \item \textbf{Memory Layer} (Wave 1B) — Milestone retrospective, friction-to-pattern registry, institutional memory
  \item \textbf{Synthesis Layer} (Wave 2) — Integration, cross-module coherence, philosophical through-line
  \item \textbf{Publication Layer} (Wave 3) — Final docs, test suite, index page
\end{enumerate}

\subsection{Deliverables}

\begin{center}
\rowcolors{2}{rowgray}{white}
\begin{tabularx}{\textwidth}{C{0.8cm}L{4cm}L{4cm}X}
\toprule
\rowcolor{primary}
\tableheader{\#} & \tableheader{Deliverable} & \tableheader{Acceptance Criteria} & \tableheader{Spec Location} \\
\midrule
D1 & Complex Plane coordinate system & $(θ, r)$ for any concept, deterministic within $\pm 0.01$ rad & Module 1 \\
D2 & `.college/` directory structure & Compiles; Mathematics seed loads without errors & Module 2 \\
D3 & Rosetta panel routing & 20 test concepts routed correctly & Module 2 \\
D4 & Chord detection algorithm & $\geq 90\%$ precision on 30 test pairs & Module 3 \\
D5 & Versine gap analysis report & Top 10 gaps verified by human review & Module 3 \\
D6 & L-system production rules & Mathematics seed expands 3 generations coherently & Module 4 \\
D7 & Angular velocity constraints & Bounded expansion, $> 95\%$ coherence & Module 4 \\
D8 & Retrospective document & All 65+ milestones with friction point + pattern & Module 5 \\
D9 & Friction-to-pattern registry & $\geq 15$ named patterns with formal descriptions & Module 5 \\
D10 & Philosophical synthesis doc & All authored works through-line mapped to architecture & Module 6 \\
D11 & Test suite & All 12 success criteria covered, 2--4 tests each & Wave 3 \\
D12 & Complete College document & Self-contained; new engineer can build from doc alone & Wave 3 \\
\bottomrule
\end{tabularx}
\end{center}

\subsection{Component Breakdown}

\begin{center}
\rowcolors{2}{rowgray}{white}
\begin{tabularx}{\textwidth}{L{4cm}L{3cm}C{2cm}C{2.5cm}}
\toprule
\rowcolor{primary}
\tableheader{Component} & \tableheader{Dependencies} & \tableheader{Model} & \tableheader{Est. Tokens} \\
\midrule
W0: Concept metadata schema & None & Haiku & 8k \\
W0: Test fixture generation & None & Haiku & 6k \\
W1A: Plane coordinate impl. & W0 schemas & Sonnet & 22k \\
W1A: College Architecture spec & W0 schemas & Sonnet & 28k \\
W1B: Milestone retrospective & W0 schemas & Opus & 35k \\
W1B: Philosophical synthesis & The Space Between & Opus & 30k \\
W1C: Chord detection algorithm & W1A Plane & Sonnet & 20k \\
W1C: Versine gap analysis & W1A Plane & Sonnet & 18k \\
W1C: L-system production rules & W1A Plane & Sonnet & 22k \\
W1C: Angular velocity constraints & W1C L-system & Sonnet & 15k \\
W2: Cross-module integration & All W1 & Opus & 40k \\
W3: Test suite compilation & W2 & Sonnet & 20k \\
W3: Final document assembly & W2 & Sonnet & 25k \\
W3: Index page & W3 docs & Sonnet & 8k \\
\midrule
\multicolumn{3}{r}{\textbf{Total estimated tokens:}} & \textbf{$\sim$297k} \\
\bottomrule
\end{tabularx}
\end{center}

\subsection{Model Assignment Rationale}

\begin{center}
\rowcolors{2}{rowgray}{white}
\begin{tabularx}{\textwidth}{L{2cm}C{1.5cm}L{4cm}X}
\toprule
\rowcolor{primary}
\tableheader{Model} & \tableheader{Share} & \tableheader{Assigned To} & \tableheader{Rationale} \\
\midrule
Opus & $\sim$32\% & Retrospective, philosophical synthesis, cross-module integration & High judgment: surfacing patterns from 65 milestones requires architectural intuition; philosophical through-line requires coherence across complex abstractions \\
Sonnet & $\sim$58\% & Plane implementation, College spec, chord/versine/L-system algorithms, test suite, final assembly & Structured content with clear specifications; document assembly; algorithmic implementations \\
Haiku & $\sim$10\% & Schemas, test fixtures, scaffolding & Minimal-footprint generation; token efficiency for repeatable low-judgment work \\
\bottomrule
\end{tabularx}
\end{center}

\section{Wave Execution Plan}

\subsection{Wave Summary}

\begin{center}
\rowcolors{2}{rowgray}{white}
\begin{tabularx}{\textwidth}{C{1.2cm}L{3cm}C{2cm}L{2.5cm}X}
\toprule
\rowcolor{primary}
\tableheader{Wave} & \tableheader{Name} & \tableheader{Mode} & \tableheader{Model(s)} & \tableheader{Gate Condition} \\
\midrule
0 & Foundation & Sequential & Haiku & Schemas validated \\
1A & Architecture track & Parallel & Sonnet & Plane + College spec complete \\
1B & Memory track & Parallel & Opus & Retrospective + Synthesis complete \\
1C & Engine track & Parallel & Sonnet & Linking + Growth engines complete \\
2 & Synthesis & Sequential & Opus & All Wave 1 tracks PASS \\
3 & Publication & Sequential & Sonnet & Integration PASS; tests green \\
\bottomrule
\end{tabularx}
\end{center}

\subsection{Wave 0: Foundation}

\begin{center}
\rowcolors{2}{rowgray}{white}
\begin{tabularx}{\textwidth}{L{3cm}L{4cm}X}
\toprule
\rowcolor{secondary}
\tableheader{Task} & \tableheader{Output} & \tableheader{Notes} \\
\midrule
W0-1: Concept metadata schema & \texttt{RosettaConcept.ts} type definition & theta, r, panels, dependencies, dept \\
W0-2: Plane coordinate schema & \texttt{PlanePosition.ts} & $(\theta, r)$, quadrant, angular velocity \\
W0-3: Deep link schema & \texttt{ConceptLink.ts} & Source, target, chord length, link type \\
W0-4: Test fixture set & \texttt{fixtures/concepts.json} & 30 labeled concept pairs for chord testing \\
W0-5: L-system schema & \texttt{LSystemRule.ts} & Axiom, production rules, generation metadata \\
\bottomrule
\end{tabularx}
\end{center}

\subsection{Wave 1A: Architecture Track (Parallel)}

\begin{center}
\rowcolors{2}{rowgray}{white}
\begin{tabularx}{\textwidth}{L{3cm}L{3cm}X}
\toprule
\rowcolor{secondary}
\tableheader{Task} & \tableheader{Output} & \tableheader{Notes} \\
\midrule
W1A-1: Plane coordinate impl. & \texttt{src/college/plane.ts} & $(θ,r)$ assignment for any concept input \\
W1A-2: Euler composition & \texttt{src/college/compose.ts} & Skill combination via complex multiplication \\
W1A-3: `.college/` directory & \texttt{.college/} directory with schema & Math, Cooking dept stubs \\
W1A-4: Math Dept seed & \texttt{.college/math/seed.yaml} & Full Mathematics seed (7 panels) \\
W1A-5: Rosetta routing & \texttt{src/college/router.ts} & Maps plane position to panel selection \\
\bottomrule
\end{tabularx}
\end{center}

\subsection{Wave 1B: Memory Track (Parallel, Opus)}

\begin{center}
\rowcolors{2}{rowgray}{white}
\begin{tabularx}{\textwidth}{L{3cm}L{3cm}X}
\toprule
\rowcolor{secondary}
\tableheader{Task} & \tableheader{Output} & \tableheader{Notes} \\
\midrule
W1B-1: Milestone audit (0--32) & \texttt{docs/retro/milestones-0-32.md} & First 33 milestones, friction + pattern \\
W1B-2: Milestone audit (33--65) & \texttt{docs/retro/milestones-33-65.md} & Milestones 33--65+, friction + pattern \\
W1B-3: Friction-to-pattern registry & \texttt{docs/retro/patterns.md} & $\geq$15 named patterns, formal descriptions \\
W1B-4: Philosophical synthesis & \texttt{docs/philosophy/through-line.md} & All 5 authored works mapped to architecture \\
\bottomrule
\end{tabularx}
\end{center}

\subsection{Wave 1C: Engine Track (Parallel)}

\begin{center}
\rowcolors{2}{rowgray}{white}
\begin{tabularx}{\textwidth}{L{3cm}L{3cm}X}
\toprule
\rowcolor{secondary}
\tableheader{Task} & \tableheader{Output} & \tableheader{Notes} \\
\midrule
W1C-1: Chord detection & \texttt{src/college/chord.ts} & Chord/arc ratio, deep link threshold \\
W1C-2: Versine gap analysis & \texttt{src/college/versine.ts} & Missing bridge detection, ranked output \\
W1C-3: Cross-dept pointer format & \texttt{src/college/links.ts} & Link type: chord, functor, exsecant \\
W1C-4: L-system rules (Math) & \texttt{src/college/lsystem.ts} & Production rules for Math Department \\
W1C-5: Angular velocity constraints & \texttt{src/college/growth.ts} & Bounded expansion, Mandelbrot gate \\
W1C-6: Calibration loop integration & \texttt{src/college/calibrate.ts} & Observe→Compare→Adjust→Record \\
\bottomrule
\end{tabularx}
\end{center}

\subsection{Wave 2: Synthesis}

\begin{center}
\rowcolors{2}{rowgray}{white}
\begin{tabularx}{\textwidth}{L{3cm}L{3cm}X}
\toprule
\rowcolor{secondary}
\tableheader{Task} & \tableheader{Output} & \tableheader{Notes} \\
\midrule
W2-1: Plane $\leftrightarrow$ routing integration & Integration test suite & Verify plane positions drive panel routing \\
W2-2: Chord $\leftrightarrow$ seed integration & Cross-dept expansion test & Seed growth respects chord boundaries \\
W2-3: Retro $\leftrightarrow$ architecture mapping & Architecture decision log & Friction points directly reference code \\
W2-4: Synthesis document & \texttt{docs/synthesis.md} & All modules coherent; through-line verified \\
\bottomrule
\end{tabularx}
\end{center}

\subsection{Wave 3: Publication}

\begin{center}
\rowcolors{2}{rowgray}{white}
\begin{tabularx}{\textwidth}{L{3cm}L{3cm}X}
\toprule
\rowcolor{secondary}
\tableheader{Task} & \tableheader{Output} & \tableheader{Notes} \\
\midrule
W3-1: Test suite compilation & \texttt{tests/college/} & Full 12-criterion coverage \\
W3-2: Final College document & \texttt{docs/college-complete.md} & Self-contained; new-engineer-ready \\
W3-3: README update & \texttt{README.md} patch & College Structure section added \\
W3-4: Index page & \texttt{docs/college/index.html} & Navigable entry point \\
\bottomrule
\end{tabularx}
\end{center}

\subsection{Cache Optimization Strategy}

\begin{itemize}[noitemsep]
  \item Wave 0 outputs (schemas) loaded once and cached across all Wave 1 tracks. Do not re-derive.
  \item \textit{The Space Between} plane definitions are static reference; cache for entire session.
  \item Rosetta panel routing table (W1A-5) cached before Wave 1C begins; chord calculations reference it.
  \item Retrospective audit (W1B-1, W1B-2) split into two sub-sessions with explicit handoff; each sub-session receives the previous's friction-to-pattern summary as preamble.
  \item Wave 2 integration receives synthesis summary from each Wave 1 track, not full outputs. Token budget: 3k summary per track.
\end{itemize}

\subsection{Token Budget}

\begin{center}
\rowcolors{2}{rowgray}{white}
\begin{tabularx}{\textwidth}{L{3cm}C{2.5cm}C{2.5cm}X}
\toprule
\rowcolor{primary}
\tableheader{Wave} & \tableheader{Input Tokens} & \tableheader{Output Tokens} & \tableheader{Notes} \\
\midrule
Wave 0 & 8k & 15k & Schemas + fixtures \\
Wave 1A & 25k & 40k & Plane + College \\
Wave 1B & 60k & 55k & Retro audit (long context) \\
Wave 1C & 20k & 35k & Engine implementations \\
Wave 2 & 30k & 25k & Synthesis (summaries only) \\
Wave 3 & 25k & 30k & Final assembly \\
\midrule
\textbf{Total} & \textbf{$\sim$168k} & \textbf{$\sim$200k} & \textbf{$\sim$368k total} \\
\bottomrule
\end{tabularx}
\end{center}

\subsection{Dependency Graph}

\begin{asciibox}
  W0 (Foundation)
   ├──────────────────────────────────┐
   │                                  │
  W1A (Architecture)    W1B (Memory)  W1C (Engines)
  [Plane, College]       [Retro,Synth]  [Links, Growth]
         │                   │              │
         └───────────────────┼──────────────┘
                             │
                         W2 (Synthesis)
                             │
                         W3 (Publication)

  Legend:
  ──→  Data dependency (output of A is input of B)
  W1A, W1B, W1C run fully in parallel (no cross-deps)
  W2 requires ALL W1 tracks to PASS
\end{asciibox}

\section{Test Plan}

\subsection{Test Categories}

\begin{center}
\rowcolors{2}{rowgray}{white}
\begin{tabularx}{\textwidth}{L{3.5cm}C{1.5cm}C{2cm}X}
\toprule
\rowcolor{primary}
\tableheader{Category} & \tableheader{Count} & \tableheader{Priority} & \tableheader{Failure Action} \\
\midrule
Safety-critical & 4 & BLOCK & Mission halt; do not proceed \\
Plane geometry & 8 & HIGH & BLOCK Wave 1C if fails \\
College structure & 6 & HIGH & BLOCK Wave 2 if fails \\
Deep linking & 8 & HIGH & BLOCK Wave 2 if fails \\
Fractal growth & 6 & HIGH & BLOCK Wave 3 if fails \\
Retrospective completeness & 4 & MEDIUM & Flag; human review before Wave 3 \\
Integration (cross-module) & 8 & HIGH & Rerun Wave 2 with fixes \\
Edge cases & 10 & LOW & Document and defer \\
\midrule
\textbf{Total} & \textbf{54} & --- & --- \\
\bottomrule
\end{tabularx}
\end{center}

\subsection{Safety-Critical Tests}

\begin{center}
\rowcolors{2}{rowgray}{white}
\begin{tabularx}{\textwidth}{L{1.5cm}L{4cm}X}
\toprule
\rowcolor{primary}
\tableheader{ID} & \tableheader{Verifies} & \tableheader{Expected Outcome} \\
\midrule
SC-01 & Cultural attribution accuracy & Coast Salish references use nation-specific language; no generic ``Indigenous'' \\
SC-02 & Authored work citation integrity & All references to authored works use exact titles and pen name \\
SC-03 & Learner data privacy & Calibration profiles do not leak across learner sessions; default private \\
SC-04 & Incoherence gate (Mandelbrot) & Unbounded seed expansion is terminated, not allowed to corrupt concept graph \\
\bottomrule
\end{tabularx}
\end{center}

\subsection{Verification Matrix}

\begin{center}
\rowcolors{2}{rowgray}{white}
\begin{tabularx}{\textwidth}{L{4.5cm}L{3cm}C{2cm}}
\toprule
\rowcolor{primary}
\tableheader{Success Criterion} & \tableheader{Test IDs} & \tableheader{Status} \\
\midrule
Plane coordinate system deterministic & PG-01, PG-02, INT-01 & Pending \\
College seed loads without errors & CS-01, CS-02 & Pending \\
Rosetta routing correct & CS-03, CS-04, INT-02 & Pending \\
Chord detection $\geq 90\%$ precision & DL-01, DL-02, DL-03 & Pending \\
Versine gap top 10 verified & DL-04, DL-05 & Pending \\
L-system expands 3 generations & FG-01, FG-02, FG-03 & Pending \\
Angular velocity constraint holds & FG-04, FG-05, SC-04 & Pending \\
Retrospective document complete & RC-01, RC-02 & Pending \\
Friction-to-pattern registry $\geq 15$ & RC-03, RC-04 & Pending \\
Authored works through-line mapped & INT-07, INT-08 & Pending \\
Cross-dept navigation 95\%+ correct & DL-06, DL-07, INT-04 & Pending \\
Self-contained for new engineer & INT-08, INT-09 & Pending \\
\bottomrule
\end{tabularx}
\end{center}

\section{Mission Crew Manifest}

\begin{center}
\rowcolors{2}{rowgray}{white}
\begin{tabularx}{\textwidth}{L{2cm}L{3.5cm}X}
\toprule
\rowcolor{primary}
\tableheader{Role} & \tableheader{Agent} & \tableheader{Responsibility} \\
\midrule
FLIGHT & AGNUS-ORCH & Mission direction; go/no-go at every wave boundary; architecture coherence \\
PLAN & PLAN-WAVE & Wave decomposition; parallel track optimization; cache exploitation \\
TOPO & TOPO-STRUCT & Agent team topology; mid-mission restructuring if W1 tracks diverge \\
ARCH & ARCH-PLANE & Cross-cutting decisions about Complex Plane coordinate system \\
SCOUT & SCOUT-LINK & Source gathering for retrospective; milestone history compilation \\
EXEC-A & DENISE-ARCH & Document production: Wave 1A outputs \\
EXEC-B & 68K-ENGINE & Document production: Wave 1C outputs \\
CRAFT-FRACTAL & PAULA-FRACTAL & L-system formalization; fractal depth boundary conditions \\
CRAFT-RETRO & RAVEN-RETRO & Milestone retrospective; friction-to-pattern extraction \\
CRAFT-PHILO & FOX-PHILO & Philosophical synthesis; through-line coherence across authored works \\
INTEG & INTEG-CHORD & Cross-module relationship validation; chord/plane/seed coherence \\
VERIFY & VERIFY-TEST & Source validation; test suite execution; accuracy checking \\
SECURE & SECURE-CULTURAL & Cultural attribution enforcement; SC-01, SC-02 mandatory-pass \\
ANALYST & ANALYST-PATTERN & Cross-module pattern detection in retrospective data \\
CAPCOM & CAPCOM-HUMAN & Human approval at wave boundaries; ONLY channel to Tibsfox \\
BUDGET & BUDGET-TOK & Token budget tracking; session planning; cache TTL exploitation \\
SURGEON & SURGEON-HEALTH & Research quality degradation detection; coherence monitoring \\
RETRO & RETRO-WAVE & Post-wave analysis; lessons learned within this mission \\
LOG & CHRONICLE & Audit trail; commit messages; milestone summary for v1.50 \\
\bottomrule
\end{tabularx}
\end{center}

\section{Execution Summary}

\begin{center}
\rowcolors{2}{rowgray}{white}
\begin{tabularx}{\textwidth}{L{5cm}X}
\toprule
\rowcolor{primary}
\tableheader{Metric} & \tableheader{Value} \\
\midrule
Activation profile & Fleet (6 modules, 19 roles) \\
Total deliverables & 12 (D1--D12) \\
Total tests & 54 (4 safety-critical, 50 functional/integration) \\
Estimated context windows & 18--22 across 6--8 sessions \\
Estimated token budget & $\sim$368k total \\
Parallel tracks & 3 (Wave 1A, 1B, 1C simultaneous) \\
Hard dependencies & W0 $\to$ W1; W1 complete $\to$ W2; W2 $\to$ W3 \\
Safety-critical blocks & 4 (SC-01 through SC-04, all mandatory-pass) \\
Opus allocation & $\sim$32\% (synthesis, retrospective, philosophical) \\
Sonnet allocation & $\sim$58\% (implementations, documents, assembly) \\
Haiku allocation & $\sim$10\% (schemas, scaffolding, fixtures) \\
Milestone designation & v1.50 Half B, Chapter: College Foundations \\
GSD repo branch & dev (\texttt{github.com/Tibsfox/gsd-skill-creator}) \\
\bottomrule
\end{tabularx}
\end{center}

\vspace{2em}
\begin{quotebox}
\textit{``The universe is an L-system that accidentally grew consciousness. GSD is a grammar for
consciousness to intentionally grow what comes next.''}

\vspace{0.5em}
--- \textit{The Space Between}, closing statement

\vspace{1em}
The College of Knowledge is not a curriculum. It is a grammar. It does not tell a learner what
to think. It provides the production rules --- the Rosetta panels, the calibration loops, the
angular velocity constraints --- that allow the learner's own derivation to grow something the
universe has never seen before. Every derivation is unique. Every tree is different. The grammar
is the same.

The space between departments is not a gap. It is where the chords live. And the chords are the
curriculum.
\end{quotebox}

\end{document}
